\section{Forward-Only Attribution: Complete Results and Boundaries}
\label{app:forward-attribution}

This appendix expands the compact comparison in
Section~\ref{sec:forward-attribution}.
It records the strict common-support protocol, complete baseline results,
an official-protocol FunnyBirds reproduction check, the endpoint-versus-trajectory
ablation, full-scale ImageNet validation, measured compute, large-model scaling,
and the PartImageNet boundary case.

\subsection{Protocol and strict common support}

FunnyBirds and ImageNet-1k use the same three architectures: ResNet-50, VGG-16, and ViT-B/16. FunnyBirds evaluates semantic-part rankings against two held-out operators, Telea inpainting and background-texture replacement. ImageNet-1k IDSDS partitions every image into a fixed $4\times4$ grid and evaluates the ranking of 16 patch scores against in-domain single-deletion effects. We first filter to images classified correctly by each model. We then intersect image IDs across the methods included in the strict comparison. The resulting support contains 499, 497, and 488 FunnyBirds images, and 7,663, 7,189, and 8,285 ImageNet images, respectively.

For each image, the primary metric is Spearman correlation between feature attribution and the benchmark target. We average within each model and then macro-average equally across architectures. All intervals use 1,000 paired image-cluster bootstrap replicates. We never average the FunnyBirds and ImageNet columns into one leaderboard.

Endpoint $M_k=|d_k|$ is an endpoint-only reference. On ImageNet IDSDS, the endpoint audit gives $\max_k|d_k-g_k|=0$ and confirms that every stored endpoint effect uses the same deletion pair as the evaluation target. The direct target-derived $|g_k|$ and independently persisted DECAF $M_k$ differ only by a mean numerical discrepancy of $6.0\times10^{-6}$.

\subsection{Complete cross-dataset method comparison}
\label{app:forward-attribution-comparison}

\begin{table*}[t]
\centering
\caption{Complete strict-common-support attribution results. Entries are equal-architecture macro-average Spearman correlations with $95\%$ paired image-cluster bootstrap intervals. Endpoint $M$ is an endpoint-only reference. KernelSHAP-512 is a deletion-game reference.}
\label{tab:forward-attribution-complete}
\footnotesize
\setlength{\tabcolsep}{4.5pt}
\begin{tabular}{llcc}
\toprule
Method & Access & FunnyBirds $\rho$ [95\% CI] & ImageNet-1k IDSDS $\rho$ [95\% CI] \\
\midrule
\multicolumn{4}{l}{\emph{DECAF}} \\
\decaf-3 & Forward only & 0.372 [0.353, 0.392] & 0.359 [0.355, 0.364] \\
\decaf-5 & Forward only & 0.403 [0.385, 0.422] & 0.367 [0.362, 0.371] \\
\decaf-9 & Forward only & 0.406 [0.388, 0.425] & 0.379 [0.375, 0.383] \\
\addlinespace
\multicolumn{4}{l}{\emph{General-purpose attribution}} \\
Input$\times$Gradient & Backward & 0.019 [$-$0.002, 0.042] & 0.098 [0.094, 0.101] \\
IG-16 & Backward & 0.266 [0.246, 0.285] & 0.238 [0.235, 0.242] \\
IG-32 & Backward & 0.271 [0.251, 0.290] & 0.242 [0.238, 0.245] \\
IG-U-32 & Backward & 0.200 [0.178, 0.222] & 0.295 [0.292, 0.299] \\
DeepLIFT & Backward & 0.197 [0.176, 0.217] & 0.341 [0.338, 0.344] \\
GradientSHAP & Backward & 0.226 [0.204, 0.246] & 0.236 [0.233, 0.240] \\
SmoothGrad-16 & Backward & $-$0.045 [$-$0.065, $-$0.023] & $-$0.015 [$-$0.018, $-$0.011] \\
RISE-512 & Sampling & 0.302 [0.283, 0.321] & 0.179 [0.175, 0.183] \\
RISE-U-512 & Sampling & 0.294 [0.274, 0.315] & 0.111 [0.107, 0.115] \\
\addlinespace
\multicolumn{4}{l}{\emph{Endpoint-only reference}} \\
Endpoint $M$ & Endpoint only & 0.324 [0.303, 0.344] & 0.371 [0.366, 0.377] \\
\addlinespace
\multicolumn{4}{l}{\emph{Deletion-game reference}} \\
KernelSHAP-512 & Sampling & 0.299 [0.278, 0.319] & 0.447 [0.443, 0.451] \\
\bottomrule
\end{tabular}
\end{table*}

\paragraph{Published-benchmark sanity check.}
On the original full-scale IDSDS ResNet-50 setting, our IG and IG-U scores are 0.194 and 0.252, closely matching the published values of 0.196 and 0.255~\citep{hesse2024benchmarking}. This agreement provides an external check on our ImageNet implementation. Our FunnyBirds evaluation instead uses held-out intervention operators and strict common support, so it is not intended to reproduce the dataset's native part-based evaluation protocol~\citep{hesse2023funnybirds}. We therefore compare its baseline values only within our registered protocol.


The stronger IG baseline changes across datasets. IG-32 exceeds IG-U-32 on FunnyBirds, while IG-U-32 is stronger on IDSDS. \decaf-5 exceeds both variants on both datasets. KernelSHAP exhibits the opposite transfer pattern: it is strongest on IDSDS, whose target is the deletion game it queries, but is weaker than endpoint $M$ and DECAF under the held-out FunnyBirds operators.

\paragraph{Native-protocol reproduction check.}
The FunnyBirds column above evaluates attribution transfer to two held-out
intervention operators, rather than the dataset's native Single Deletion
protocol. We therefore ran a separate implementation check using the official
FunnyBirds setting: the released checkpoints, native semantic part removals,
and RISE with 6,000 masks, an $8\times8$ coarse grid, and mask probability
$p=0.1$~\citep{hesse2023funnybirds}. As shown in
Table~\ref{tab:funnybirds-native-reproduction}, the reproduced RISE scores
differ from the published values by less than $0.005$ for every architecture.
This check supports the implementation fidelity of the RISE baseline used in
our benchmark; it does not imply that the held-out score of $0.302$ should
numerically match the native Single Deletion scores, because the two
evaluations use different intervention targets and RISE configurations.

\begin{table}[t]
\centering
\caption{
FunnyBirds native Single Deletion reproduction and native-target audit.
\textbf{Top:} the official RISE-6000 setting reproduces the published
Single Deletion scores within $0.005$ on all three architectures.
\textbf{Bottom:} raw Spearman correlation with the same native part-removal
target. Endpoint $M$ is included to show the strong target alignment of this
evaluation. These native-target results are a reproduction and boundary check,
not a replacement for the held-out FunnyBirds comparison in
Table~\ref{tab:forward-attribution-main}.
}
\label{tab:funnybirds-native-reproduction}
\small
\setlength{\tabcolsep}{4.5pt}

\begin{tabular}{lccc}
\toprule
Architecture & Published RISE SD & Reproduced RISE SD & $|\Delta|$ \\
\midrule
ResNet-50 & 0.560 & 0.558 & 0.002 \\
VGG16     & 0.730 & 0.726 & 0.004 \\
ViT-B/16  & 0.790 & 0.788 & 0.002 \\
\midrule
\multicolumn{4}{c}{\textit{Native-target raw Spearman $\rho$}} \\
\midrule
Method & ResNet-50 & VGG16 & ViT-B/16 \\
\midrule
Endpoint $M$ & 0.757 & 1.000 & 0.983 \\
\decaf-3     & 0.484 & 0.926 & 0.909 \\
\decaf-5     & 0.565 & 0.898 & 0.864 \\
\decaf-9     & 0.616 & 0.901 & 0.850 \\
RISE-6000    & 0.115 & 0.452 & 0.576 \\
\bottomrule
\end{tabular}
\end{table}

The native and held-out FunnyBirds experiments answer different questions.
The native target is defined by the same semantic part-removal contrast used
to construct the endpoint, whereas the main benchmark evaluates transfer to
held-out inpainting and texture-replacement interventions. The native audit
therefore verifies baseline implementation and exposes the target-aligned
boundary; the held-out benchmark remains the test of attribution transfer.

\subsection{Endpoint versus trajectory}
\label{app:endpoint-versus-trajectory}



\begin{table}[t]
\centering
\caption{Paired endpoint-versus-trajectory tests. Differences are DECAF minus endpoint $M$. Positive values indicate ordinary-attribution information beyond the endpoint.}
\label{tab:endpoint-trajectory-paired}
\footnotesize
\setlength{\tabcolsep}{4pt}
\begin{tabular}{llrrr}
\toprule
Dataset & Contrast & Mean $\Delta$ & 95\% CI & Model wins \\
\midrule
FunnyBirds & \decaf-3 $-M$ & +0.0488 & [0.0334, 0.0642] & 2/3 \\
FunnyBirds & \decaf-5 $-M$ & +0.0797 & [0.0637, 0.0962] & 2/3 \\
FunnyBirds & \decaf-9 $-M$ & +0.0828 & [0.0671, 0.0993] & 2/3 \\
\addlinespace
ImageNet-1k & \decaf-3 $-M$ & $-$0.0119 & [$-$0.0142, $-$0.0096] & 1/3 \\
ImageNet-1k & \decaf-5 $-M$ & $-$0.0045 & [$-$0.0072, $-$0.0017] & 1/3 \\
ImageNet-1k & \decaf-9 $-M$ & +0.0073 & [0.0045, 0.0101] & 2/3 \\
\bottomrule
\end{tabular}
\end{table}

FunnyBirds is the clean trajectory-value test because its evaluation operators differ from the explanation endpoint. All three DECAF grids significantly exceed $M$. Five stages capture $96.3\%$ of the nine-stage gain over the endpoint. IDSDS evaluates the same deletion contrast that defines $d_k$. Endpoint $M$ is therefore already highly aligned with its target. Three and five stages preserve most of that signal, while nine stages add a small but stable gain.

\begin{table*}[t]
\centering
\caption{Endpoint-versus-trajectory ablation by architecture. Best $\Delta$ is the largest DECAF-3/5/9 Spearman minus endpoint $M$.}
\label{tab:endpoint-trajectory-per-model}
\footnotesize
\setlength{\tabcolsep}{4.2pt}
\begin{tabular}{llrrrrr}
\toprule
Dataset & Model & Endpoint $M$ & \decaf-3 & \decaf-5 & \decaf-9 & Best $\Delta$ \\
\midrule
FunnyBirds & ResNet-50 & 0.388 & 0.425 & 0.447 & 0.451 & +0.064 \\
FunnyBirds & VGG-16 & 0.306 & 0.283 & 0.298 & 0.293 & $-$0.008 \\
FunnyBirds & ViT-B/16 & 0.277 & 0.409 & 0.465 & 0.475 & +0.197 \\
ImageNet-1k & ResNet-50 & 0.384 & 0.362 & 0.370 & 0.387 & +0.002 \\
ImageNet-1k & VGG-16 & 0.671 & 0.640 & 0.640 & 0.644 & $-$0.027 \\
ImageNet-1k & ViT-B/16 & 0.058 & 0.077 & 0.090 & 0.105 & +0.047 \\
\bottomrule
\end{tabular}
\end{table*}

The trajectory gain is architecture-dependent in the tested models. It is largest for ViT-B/16 on both datasets, modest for ResNet-50, and negative for VGG-16. This pattern is empirical rather than a general architecture law, but it shows that the value of intermediate responses is not uniform across model classes.

The separate native FunnyBirds audit in
Appendix~\ref{app:forward-attribution-comparison} provides the complementary
endpoint-aligned case: when the evaluation target is the exact native
part-removal contrast, endpoint $M$ reaches $\rho=0.913$ and exceeds all
trajectory summaries.

Together, the two FunnyBirds evaluations isolate the role of the trajectory:
the endpoint best recovers its own intervention effect, while intermediate
responses add value when attribution must transfer to a different intervention.

\subsection{Full-scale ImageNet validation}

\begin{table}[t]
\centering
\caption{Full 50,000-image ImageNet validation scale check. Values are equal-architecture macro-average IDSDS over 115,876 correctly classified model--image units.}
\label{tab:idsds-full50k}
\footnotesize
\setlength{\tabcolsep}{4pt}
\begin{tabular}{lrrrr}
\toprule
Method & Macro $\rho$ [95\% CI] & ResNet-50 & VGG-16 & ViT-B/16 \\
\midrule
Endpoint $M$ & 0.3708 [0.3685, 0.3731] & 0.3778 & 0.6731 & 0.0615 \\
\decaf-5 & 0.3633 [0.3613, 0.3653] & 0.3609 & 0.6395 & 0.0894 \\
IG-U-32 & 0.2942 [0.2926, 0.2958] & 0.2522 & 0.4939 & 0.1366 \\
IG-32 & 0.2397 [0.2380, 0.2411] & 0.1941 & 0.4002 & 0.1247 \\
\bottomrule
\end{tabular}
\end{table}

The 10,000-image conclusions are not a favorable-subset artifact. The full validation set preserves the ordering $M > \decaf\text{-}5 > \mathrm{IG\text{-}U\text{-}32} > \mathrm{IG\text{-}32}$. The full-scale paired difference between \decaf-5 and $M$ is $-0.0075$ with a $95\%$ interval of $[-0.0086,-0.0062]$.

\subsection{Measured compute and large-model scaling}

\begin{table*}[!tb]
\centering
\caption{Measured ImageNet-1k compute. Rows per image count model-forward input rows after batching. Timing is compute-only and macro-averaged across ResNet-50, VGG-16, and ViT-B/16.}
\label{tab:forward-attribution-compute}
\footnotesize
\setlength{\tabcolsep}{4.5pt}
\begin{tabular}{llrrrr}
\toprule
Method & Access & Rows/image & Backward? & ms/image $\downarrow$ & Peak GiB $\downarrow$ \\
\midrule
\decaf-3 & Forward only & 51 & No & 22.1 & 11.2 \\
\decaf-5 & Forward only & 85 & No & 36.7 & 11.2 \\
\decaf-9 & Forward only & 153 & No & 65.9 & 11.2 \\
Input$\times$Gradient & Backward & 1 & Yes & 1.4 & 2.3 \\
DeepLIFT & Backward & 2 & Yes & 3.1 & 6.6 \\
GradientSHAP & Backward & 16 & Yes & 14.4 & 25.9 \\
IG-32 & Backward & 32 & Yes & 28.7 & 50.4 \\
IG-U-32 & Backward & 32 & Yes & 28.7 & 50.4 \\
SmoothGrad-16 & Backward & 16 & Yes & 16.1 & 25.9 \\
RISE-512 & Sampling & 512 & No & 216.6 & 14.4 \\
KernelSHAP-512 & Sampling & 512 & No & 216.4 & 14.4 \\
\bottomrule
\end{tabular}
\end{table*}

\decaf-3 is the speed-oriented Pareto point. Relative to IG-32, it has higher quality on both datasets, is $1.30\times$ faster, and uses $4.50\times$ less peak memory. \decaf-5 trades additional latency for higher quality while keeping the same inference-scale memory. KernelSHAP-512 is $5.89\times$ slower than \decaf-5 and uses six times as many forward rows.




\subsection{PartImageNet as a task-aligned boundary case}

PartImageNet supplies semantic part masks and evaluates held-out part-removal effects. Direct part removal and coalition methods are therefore unusually aligned with the evaluation target. We retain the benchmark as a boundary case rather than using it to define the main general-purpose comparison.

\begin{table}[!tb]
\centering
\caption{PartImageNet strict-common-support Spearman. Part-removal and coalition methods exploit the supplied semantic part groups and are directly aligned with the held-out part-removal target.}
\label{tab:partimagenet-boundary}
\footnotesize
\setlength{\tabcolsep}{4pt}
\begin{tabular}{lr}
\toprule
Method & Spearman [95\% CI] \\
\midrule
Part-LIME-1000 & 0.478 [0.458, 0.497] \\
Part Occlusion & 0.477 [0.458, 0.496] \\
Exact Part-Shapley & 0.433 [0.412, 0.454] \\
KernelSHAP-512 & 0.426 [0.405, 0.447] \\
Endpoint $M$ & 0.364 [0.342, 0.385] \\
\decaf-9 & 0.363 [0.341, 0.384] \\
\decaf-5 & 0.358 [0.337, 0.379] \\
\decaf-3 & 0.350 [0.329, 0.370] \\
RISE-512 & 0.299 [0.276, 0.320] \\
IG-32 & 0.289 [0.270, 0.311] \\
\bottomrule
\end{tabular}
\end{table}

This result marks the method boundary clearly. When perfect semantic parts are already supplied and the target is itself a part-removal effect, direct part interventions can be stronger. Even in that setting, the DECAF trajectories remain above standard gradient-path attribution and preserve their forward-only memory advantage.
